\documentclass[11pt, a4paper]{article}

% ── Geometry & fonts ──────────────────────────────────────────────────────────
\usepackage[a4paper, top=2.5cm, bottom=2.5cm, left=2.8cm, right=2.8cm]{geometry}
\usepackage{mathpazo}          % Palatino text + Pazo math
\usepackage[T1]{fontenc}
\usepackage[utf8]{inputenc}
\usepackage{microtype}

% ── Math ──────────────────────────────────────────────────────────────────────
\usepackage{amsmath, amssymb, amsthm, mathtools}
\usepackage{bm}                % bold math
\usepackage{bbm}               % indicator 1

% ── Colours & boxes ───────────────────────────────────────────────────────────
\usepackage[dvipsnames, svgnames, x11names]{xcolor}
\usepackage[most]{tcolorbox}
\usepackage{mdframed}

% ── Layout ────────────────────────────────────────────────────────────────────
\usepackage{booktabs}
\usepackage{enumitem}
\usepackage{titlesec}
\usepackage{fancyhdr}
\usepackage{abstract}
\usepackage{hyperref}
\usepackage{cleveref}

% ─────────────────────────────────────────────────────────────────────────────
% Colour palette
% ─────────────────────────────────────────────────────────────────────────────
\definecolor{NavyProof}   {HTML}{1A2E6B}
\definecolor{SteelBlue}   {HTML}{2C6FAC}
\definecolor{LightSteel}  {HTML}{DCE8F5}
\definecolor{ThmGreen}    {HTML}{2D6A1F}
\definecolor{ThmGreenBg}  {HTML}{EAF3DE}
\definecolor{LemmaAmber}  {HTML}{7A4800}
\definecolor{LemmaAmberBg}{HTML}{FEF3DC}
\definecolor{CorOrange}   {HTML}{8B3A00}
\definecolor{CorOrangeBg} {HTML}{FDEEE0}
\definecolor{PropPink}    {HTML}{7A1140}
\definecolor{PropPinkBg}  {HTML}{FBEAF0}
\definecolor{RemarkGray}  {HTML}{4A4845}
\definecolor{RemarkGrayBg}{HTML}{F3F2EE}
\definecolor{DefnPurple}  {HTML}{3B1FA8}
\definecolor{DefnPurpleBg}{HTML}{EEECFD}
\definecolor{RuleLine}    {HTML}{CCCCCC}

% ─────────────────────────────────────────────────────────────────────────────
% Theorem environments via tcolorbox
% ─────────────────────────────────────────────────────────────────────────────
\tcbuselibrary{theorems, skins, breakable}

\newcounter{mainthm}[section]
\renewcommand{\themainthm}{\thesection.\arabic{mainthm}}

% Theorem
\newtcbtheorem[use counter=mainthm, number within=section]
  {theorem}{Theorem}{
    enhanced, breakable,
    colback=LightSteel, colframe=SteelBlue,
    fonttitle=\bfseries\color{NavyProof},
    coltitle=NavyProof,
    attach boxed title to top left={yshift=-2mm, xshift=6mm},
    boxed title style={colback=SteelBlue, colframe=SteelBlue, sharp corners},
    sharp corners=south, arc=4pt,
    left=6pt, right=6pt, top=8pt, bottom=6pt,
    before upper={\setlength{\parindent}{0pt}\setlength{\parskip}{4pt}},
  }{thm}

% Lemma
\newtcbtheorem[use counter=mainthm, number within=section]
  {lemma}{Lemma}{
    enhanced, breakable,
    colback=ThmGreenBg, colframe=ThmGreen,
    fonttitle=\bfseries\color{ThmGreen},
    coltitle=ThmGreen,
    attach boxed title to top left={yshift=-2mm, xshift=6mm},
    boxed title style={colback=ThmGreen, colframe=ThmGreen, sharp corners},
    sharp corners=south, arc=4pt,
    left=6pt, right=6pt, top=8pt, bottom=6pt,
    before upper={\setlength{\parindent}{0pt}\setlength{\parskip}{4pt}},
  }{lem}

% Proposition
\newtcbtheorem[use counter=mainthm, number within=section]
  {proposition}{Proposition}{
    enhanced, breakable,
    colback=PropPinkBg, colframe=PropPink,
    fonttitle=\bfseries\color{PropPink},
    coltitle=PropPink,
    attach boxed title to top left={yshift=-2mm, xshift=6mm},
    boxed title style={colback=PropPink, colframe=PropPink, sharp corners},
    sharp corners=south, arc=4pt,
    left=6pt, right=6pt, top=8pt, bottom=6pt,
    before upper={\setlength{\parindent}{0pt}\setlength{\parskip}{4pt}},
  }{prop}

% Corollary
\newtcbtheorem[use counter=mainthm, number within=section]
  {corollary}{Corollary}{
    enhanced, breakable,
    colback=CorOrangeBg, colframe=CorOrange,
    fonttitle=\bfseries\color{CorOrange},
    coltitle=CorOrange,
    attach boxed title to top left={yshift=-2mm, xshift=6mm},
    boxed title style={colback=CorOrange, colframe=CorOrange, sharp corners},
    sharp corners=south, arc=4pt,
    left=6pt, right=6pt, top=8pt, bottom=6pt,
    before upper={\setlength{\parindent}{0pt}\setlength{\parskip}{4pt}},
  }{cor}

% Definition box
\newtcbtheorem[use counter=mainthm, number within=section]
  {definition}{Definition}{
    enhanced, breakable,
    colback=DefnPurpleBg, colframe=DefnPurple,
    fonttitle=\bfseries\color{DefnPurple},
    coltitle=DefnPurple,
    attach boxed title to top left={yshift=-2mm, xshift=6mm},
    boxed title style={colback=DefnPurple, colframe=DefnPurple, sharp corners},
    sharp corners=south, arc=4pt,
    left=6pt, right=6pt, top=8pt, bottom=6pt,
    before upper={\setlength{\parindent}{0pt}\setlength{\parskip}{4pt}},
  }{def}

% Remark (plain gray)
\newtcbtheorem[use counter=mainthm, number within=section]
  {remark}{Remark}{
    enhanced, breakable,
    colback=RemarkGrayBg, colframe=RemarkGray,
    fonttitle=\bfseries\itshape\color{RemarkGray},
    coltitle=RemarkGray,
    attach boxed title to top left={yshift=-2mm, xshift=6mm},
    boxed title style={colback=RemarkGray, colframe=RemarkGray, sharp corners},
    sharp corners=south, arc=4pt,
    left=6pt, right=6pt, top=8pt, bottom=6pt,
    before upper={\setlength{\parindent}{0pt}\setlength{\parskip}{4pt}},
  }{rem}

% Proof environment (left-bar style)
\tcolorboxenvironment{proof}{
  blanker, breakable,
  left=4mm,
  borderline west={2pt}{0pt}{RemarkGray},
  before skip=6pt, after skip=10pt,
}

% ─────────────────────────────────────────────────────────────────────────────
% Section styling
% ─────────────────────────────────────────────────────────────────────────────
\titleformat{\section}
  {\Large\bfseries\color{NavyProof}}
  {\thesection}{1em}{}
  [\vspace{-2pt}\textcolor{RuleLine}{\titlerule[0.6pt]}]

\titleformat{\subsection}
  {\large\bfseries\color{SteelBlue}}
  {\thesubsection}{1em}{}

\titleformat{\subsubsection}
  {\normalsize\bfseries\color{RemarkGray}}
  {\thesubsubsection}{1em}{}

% ─────────────────────────────────────────────────────────────────────────────
% Header / footer
% ─────────────────────────────────────────────────────────────────────────────
\pagestyle{fancy}
\fancyhf{}
\fancyhead[L]{\small\color{RemarkGray}\textit{HNMC-MELD: Mathematical Proof}}
\fancyhead[R]{\small\color{RemarkGray}\textit{MERC Architecture}}
\fancyfoot[C]{\small\color{RemarkGray}\thepage}
\renewcommand{\headrulewidth}{0.4pt}
\renewcommand{\headrule}{\hbox to\headwidth{\color{RuleLine}\leaders\hrule height \headrulewidth\hfill}}

% ─────────────────────────────────────────────────────────────────────────────
% Math operators & notation shorthands
% ─────────────────────────────────────────────────────────────────────────────
\DeclareMathOperator{\softmax}{softmax}
\DeclareMathOperator{\MLP}{MLP}
\DeclareMathOperator{\krank}{krank}
\DeclareMathOperator{\diag}{diag}
\DeclareMathOperator*{\argmax}{arg\,max}

\newcommand{\R}{\mathbb{R}}
\newcommand{\E}{\mathbb{E}}
\newcommand{\Prob}{\mathbb{P}}
\newcommand{\N}{\mathcal{N}}
\newcommand{\cS}{\mathcal{S}}
\newcommand{\cE}{\mathcal{E}}
\newcommand{\cL}{\mathcal{L}}
\newcommand{\cQ}{\mathcal{Q}}
\newcommand{\cH}{\mathcal{H}}
\newcommand{\lam}{\lambda}
\newcommand{\psi}{\psi}
\newcommand{\bh}{\mathbf{h}}
\newcommand{\bc}{\mathbf{c}}
\newcommand{\by}{\mathbf{y}}
\newcommand{\bz}{\mathbf{z}}
\newcommand{\bx}{\mathbf{x}}
\newcommand{\hbar}{\bar{h}}
\newcommand{\Abar}{\bar{A}}
\newcommand{\Bbar}{\bar{B}}
\newcommand{\norm}[1]{\left\lVert#1\right\rVert}
\newcommand{\abs}[1]{\left\lvert#1\right\rvert}
\newcommand{\given}{\,\middle|\,}
\newcommand{\pistar}{\pi^{*}}
\newcommand{\gamk}{\gamma_k(t)}
\newcommand{\QED}{\hfill$\blacksquare$}

% ─────────────────────────────────────────────────────────────────────────────
% Hyperref setup
% ─────────────────────────────────────────────────────────────────────────────
\hypersetup{
  colorlinks=true,
  linkcolor=SteelBlue,
  citecolor=ThmGreen,
  urlcolor=SteelBlue,
  pdftitle={HNMC-MELD Mathematical Proof},
  pdfauthor={MERC Architecture},
}

% ═════════════════════════════════════════════════════════════════════════════
\begin{document}
% ═════════════════════════════════════════════════════════════════════════════

% ── Title ────────────────────────────────────────────────────────────────────
\begin{center}
  {\LARGE\bfseries\color{NavyProof}
    Mathematical Proof of Correctness}\\[6pt]
  {\Large\color{SteelBlue}
    Hidden Neural Markov Chain with Mamba Experts}\\[4pt]
  {\large\itshape\color{RemarkGray}
    for Multimodal Emotion Recognition on MELD}\\[10pt]
  {\normalsize\color{RemarkGray}
    MERC Architecture\ $\cdot$\ 7-Class Emotion Classification}\\[6pt]
  \textcolor{RuleLine}{\rule{\linewidth}{0.8pt}}
\end{center}

% ── Abstract ─────────────────────────────────────────────────────────────────
\begin{abstract}
We present a rigorous mathematical framework establishing the theoretical
foundations of the \textbf{Hidden Neural Markov Chain} (HNMC) architecture,
wherein each hidden state is instantiated as a Mamba selective state-space
model (SSM) receiving the full trimodal context (text, audio, video) of each
utterance.  We prove seven properties:
(1)~the model defines a valid generative process over emotion label sequences;
(2)~exact marginal likelihoods are computable in $\mathcal{O}(K^2 T)$ time via
     the forward algorithm;
(3)~each Mamba expert is a universal approximator of within-regime emotion
     dynamics;
(4)~model parameters are generically identifiable up to expert-label
     permutation;
(5)~hybrid EM-backprop training monotonically improves the objective and
     converges to a stationary point;
(6)~HNMC strictly subsumes both classical HMMs and standard Mixture-of-Experts;
(7)~the Markov structure provides automatic structural regularisation of
     class-imbalanced rare emotions and monotone information gain from trimodal
     fusion.
Collectively these results establish that HNMC is theoretically well-founded
and provably more expressive than its constituent components applied
independently.
\end{abstract}

\tableofcontents
\vspace{1em}
\textcolor{RuleLine}{\rule{\linewidth}{0.4pt}}

% ═════════════════════════════════════════════════════════════════════════════
\section{Notation and Preliminaries}
% ═════════════════════════════════════════════════════════════════════════════

Let $\mathcal{U} = \{u_1,\ldots,u_T\}$ denote a MELD conversation of $T$
utterances.  Each utterance $u_t$ carries three modality streams, independently
encoded then fused:
%
\begin{align}
  \mathbf{z}_t^T &= \mathrm{Encoder}_T(u_t) \in \R^{d_T},
    \quad \text{(text: RoBERTa / token embeddings)} \notag\\
  \mathbf{z}_t^A &= \mathrm{Encoder}_A(u_t) \in \R^{d_A},
    \quad \text{(audio: wav2vec / MFCC)} \notag\\
  \mathbf{z}_t^V &= \mathrm{Encoder}_V(u_t) \in \R^{d_V},
    \quad \text{(video: FaceNet / ViT)} \notag\\
  \mathbf{c}_t   &= \MLP\!\bigl([\mathbf{z}_t^T \;\|\; \mathbf{z}_t^A
                     \;\|\; \mathbf{z}_t^V]\bigr)
                   \in \R^d.
                   \label{eq:fusion}
\end{align}
%
The MELD emotion label set is
$\mathcal{E} = \{\text{neutral},\text{joy},\text{surprise},\text{anger},
\text{sadness},\text{disgust},\text{fear}\}$, with $|\mathcal{E}|=K=7$.
We index emotions by $k\in\{1,\ldots,7\}$.

\begin{definition}{Model Parameters}{params}
The complete parameter set of HNMC is the triple
\[
  \lambda \;=\; \bigl(\pi_0,\; P,\; \{\psi_k\}_{k=1}^{K}\bigr),
\]
where:
\begin{itemize}[noitemsep]
  \item $\pi_0 \in \Delta^{K-1}$\;---\;initial hidden-state distribution
        on the $(K{-}1)$-simplex;
  \item $P \in \R^{K\times K}$,\; $P_{jk}\ge 0$,\;
        $\sum_{k=1}^{K}P_{jk}=1$\;---\;row-stochastic Markov transition
        matrix;
  \item $\psi_k = (A_k, B_k, C_k, D_k, \Delta_k)$\;---\;Mamba SSM
        parameters for expert $k$, with
        $A_k\in\R^{N\times N}$, $B_k\in\R^{N\times d}$,
        $C_k\in\R^{d\times N}$, $D_k\in\R^{d\times d}$,
        $\Delta_k>0$.
\end{itemize}
\end{definition}

\noindent
We write $\Delta^{K-1}$ for the probability simplex, $\norm{\cdot}_F$ for the
Frobenius norm, $\rho(\cdot)$ for the spectral radius, and
$\mathbbm{1}_{\{\cdot\}}$ for the indicator function.

\subsection{Mamba SSM Recurrence}

Each expert $S_k$ maintains a continuous hidden state $\bh_t^{(k)}\in\R^N$
updated via the zero-order-hold (ZOH) discretisation of a linear ODE:
%
\begin{align}
  \Abar_k({\bc}_t) &= \exp\!\bigl(\Delta_k\, A_k\bigr), \label{eq:Abar}\\
  \Bbar_k({\bc}_t) &= (\Delta_k A_k)^{-1}
                      \bigl[\exp(\Delta_k A_k)-I\bigr]\,\Delta_k B_k,
                      \label{eq:Bbar}\\
  \bh_t^{(k)}      &= \Abar_k(\bc_t)\;\bh_{t-1}^{(k)}
                      + \Bbar_k(\bc_t)\;\bc_t, \label{eq:ssm-h}\\
  \by_t^{(k)}      &= C_k\,\bh_t^{(k)} + D_k\,\bc_t \in \R^d.
                      \label{eq:ssm-y}
\end{align}
%
The key property of Mamba (S6) is that $\Abar_k$ and $\Bbar_k$ are
\emph{input-dependent}: $\Delta_k = \Delta_k(\bc_t)$ is produced by a small
projection, making the gating selective.

% ═════════════════════════════════════════════════════════════════════════════
\section{Validity as a Generative Model}
% ═════════════════════════════════════════════════════════════════════════════

\subsection{Construction of the Joint Distribution}

\begin{theorem}{Joint Distribution}{joint}
Under parameters $\lambda = (\pi_0, P, \{\psi_k\})$, the joint distribution
over the hidden-state sequence $z_{1:T}$ and context sequence $\bc_{1:T}$ is
well-defined and given by
\begin{equation}
  p(\bc_{1:T}, z_{1:T} \mid \lambda)
  \;=\;
  \pi_0(z_1)
  \prod_{t=2}^{T} P_{z_{t-1},\,z_t}
  \prod_{t=1}^{T} p\!\bigl(\bc_t \mid S_{z_t},\, \bh_{t-1}^{(z_t)}\bigr),
  \label{eq:joint}
\end{equation}
where
$p(\bc_t \mid S_k, \bh_{t-1}^{(k)}) =
 \mathcal{N}\!\bigl(\bc_t;\; W_k^\mu \bh_{t-1}^{(k)},\;
 \mathrm{diag}(\sigma_k(\bh_{t-1}^{(k)})^2)\bigr)$
is the Mamba emission likelihood for expert $k$.
\end{theorem}

\begin{proof}
We verify the three Kolmogorov conditions for a valid probability measure.

\medskip\noindent
\textbf{(i) Non-negativity.}\;
$\pi_0(z_1)\ge 0$ by definition on the simplex.
Each $P_{jk}\ge 0$ by the row-stochasticity constraint.
Each emission $p(\bc_t\mid S_k,\bh)\ge 0$ as it is a Gaussian density
(positive semi-definite covariance, positive normalising constant).

\medskip\noindent
\textbf{(ii) Normalisation.}\;
Summing over all state sequences
$z_{1:T}\in\{1,\ldots,K\}^T$:
\begin{align*}
  \sum_{z_{1:T}} p(\bc_{1:T},z_{1:T})
  &= \prod_{t=1}^{T} p(\bc_t\mid\cdots)
     \underbrace{\sum_{z_1}\pi_0(z_1)}_{=1}
     \prod_{t=2}^{T}
     \underbrace{\sum_{z_t} P_{z_{t-1},z_t}}_{=1} = \prod_{t=1}^T p(\bc_t\mid\cdots),
\end{align*}
which integrates to~$1$ over $\bc_{1:T}$ because each Gaussian emission
integrates to~$1$.

\medskip\noindent
\textbf{(iii) Marginalisation consistency.}\;
For any sub-sequence $\bc_{1:s}$ with $s<T$, marginalising over
$\bc_{s+1:T}$ and $z_{1:T}$ preserves the product structure by the
first-order Markov property: $z_t \perp z_{1:t-2} \mid z_{t-1}$.  \QED
\end{proof}

% ═════════════════════════════════════════════════════════════════════════════
\section{Correctness of the Forward Algorithm}
% ═════════════════════════════════════════════════════════════════════════════

\subsection{Forward Variable and Recursion}

\noindent
Define the \emph{forward variable}
$\gamma_k(t) \triangleq p(z_t=k,\, \bc_{1:t} \mid \lambda)$,
the joint probability of being in expert state $S_k$ at time~$t$ while having
observed the partial context sequence $\bc_{1:t}$.

\begin{theorem}{Forward Recursion Correctness}{forward}
The forward variables satisfy the initialisation
\begin{equation}
  \gamma_k(1) = \pi_0(k)\cdot p\!\bigl(\bc_1 \mid S_k, \bh_0^{(k)}\bigr),
  \label{eq:fw-init}
\end{equation}
and the recursion, for $t\ge 2$,
\begin{equation}
  \gamma_k(t)
  = p\!\bigl(\bc_t \mid S_k, \bh_{t-1}^{(k)}\bigr)
    \sum_{j=1}^{K} P_{jk}\,\gamma_j(t-1).
  \label{eq:fw-rec}
\end{equation}
The marginal likelihood is
$p(\bc_{1:T}\mid\lambda) = \sum_{k=1}^{K}\gamma_k(T)$.
The computation requires $\mathcal{O}(K^2 T)$ operations.
\end{theorem}

\begin{proof}
We proceed by induction on $t$.

\medskip\noindent
\textbf{Base case ($t=1$).}\;
$\gamma_k(1) = p(z_1=k,\bc_1)
 = p(z_1=k)\cdot p(\bc_1\mid z_1=k,\bh_0^{(k)})
 = \pi_0(k)\cdot p(\bc_1\mid S_k,\bh_0^{(k)})$. Correct.

\medskip\noindent
\textbf{Inductive step.}\;
Assume $\gamma_j(t-1)=p(z_{t-1}=j,\bc_{1:t-1})$ for all $j$.
Then:
\begin{align*}
  \gamma_k(t)
  &= p(z_t=k,\,\bc_{1:t})\\
  &= \sum_{j=1}^{K} p(z_{t-1}=j,\, z_t=k,\, \bc_{1:t})\\
  &= \sum_{j=1}^{K} p\!\bigl(\bc_t \mid z_t=k, \bh_{t-1}^{(k)}\bigr)
     \cdot p(z_t=k \mid z_{t-1}=j)
     \cdot p(z_{t-1}=j,\bc_{1:t-1})\\
  &= p\!\bigl(\bc_t \mid S_k, \bh_{t-1}^{(k)}\bigr)
     \sum_{j=1}^{K} P_{jk}\,\gamma_j(t-1).
\end{align*}
Step~3 invokes two conditional independence properties:
(a)~$\bc_t \perp \bc_{1:t-1} \mid (z_t, \bh_{t-1}^{(k)})$,
since $\bh_{t-1}^{(k)}$ is a deterministic function of $(\bc_{1:t-1},\psi_k)$
and hence contains all relevant history for expert~$k$;
(b)~$z_t \perp z_{1:t-2} \mid z_{t-1}$ (Markov property of $P$).

\medskip\noindent
\textbf{Complexity.}\;
Each $\gamma_k(t)$ requires one emission evaluation and $K$ multiply-adds.
Over $K$ states and $T$ steps: $\mathcal{O}(K^2 T)=\mathcal{O}(49T)$.
\QED
\end{proof}

% ═════════════════════════════════════════════════════════════════════════════
\section{Universal Approximation by Mamba Experts}
% ═════════════════════════════════════════════════════════════════════════════

\subsection{SSM as a Functional Approximator}

\begin{lemma}{SSM State as Sufficient History Statistic}{ssmstat}
For any expert $\psi_k$, the hidden state $\bh_t^{(k)}$ computed by
\cref{eq:ssm-h} is a deterministic function
$\bh_t^{(k)} = F_k(\bc_{1:t})$ of the entire context history $\bc_{1:t}$.
\end{lemma}

\begin{proof}
By induction: $\bh_0^{(k)}=\mathbf{0}$ (fixed).
For $t=1$: $\bh_1^{(k)} = \Bbar_k(\bc_1)\,\bc_1 = F_k(\bc_1)$.
For $t\ge2$: assuming $\bh_{t-1}^{(k)}=F_k(\bc_{1:t-1})$,
\[
  \bh_t^{(k)}
  = \Abar_k(\bc_t)\,F_k(\bc_{1:t-1}) + \Bbar_k(\bc_t)\,\bc_t
  \;\triangleq\; F_k(\bc_{1:t}).
\]
Since $\Abar_k$ and $\Bbar_k$ are deterministic functions of $\bc_t$
(via the $\Delta_k$ projection and the matrix exponential), $F_k$ is a
well-defined deterministic map.  \QED
\end{proof}

\begin{theorem}{Universal Approximation of Emotion Dynamics}{ua}
Let $f:(\R^d)^t\to\R^d$ be any continuous causal functional mapping
context histories to target outputs (e.g.\ the ideal within-regime
emotion predictor).  For any $\varepsilon>0$ and any compact set
$\mathcal{K}\subset(\R^d)^t$, there exist Mamba parameters $\psi_k$ and
state dimension $N$ such that
\[
  \sup_{\bc_{1:t}\in\mathcal{K}}
  \norm{C_k\,\bh_t^{(k)} + D_k\,\bc_t - f(\bc_{1:t})}_2 < \varepsilon.
\]
\end{theorem}

\begin{proof}
The output map $\bc_{1:t}\mapsto C_k\bh_t^{(k)}+D_k\bc_t$ is realised by a
\emph{recurrent network with input-dependent transitions}: the matrices
$\Abar_k(\bc_t)$ and $\Bbar_k(\bc_t)$ are composed of a linear projection
(for $\Delta_k$), a matrix exponential, and matrix multiplications---all
differentiable operations that can be implemented by MLP sub-networks.

\medskip\noindent
By the universal approximation theorem for recurrent networks
\cite{Siegelmann1992,Hammer2000}:
\emph{any continuous causal functional on a compact domain can be approximated
to arbitrary precision by a recurrent network with sigmoid activations and
sufficiently large hidden dimension $N$}.

\medskip\noindent
Since the MLP sub-networks computing $\Delta_k(\bc_t)$ are themselves
universal approximators for continuous functions
\cite{Cybenko1989,Hornik1989},
and since $C_k\in\R^{d\times N}$ is a linear read-out that matches any linear
functional of $\bh_t^{(k)}$, the composition approximates $f$ on
$\mathcal{K}$ for sufficiently large $N$.
The $\mathcal{O}(T)$ complexity of the recurrence is preserved at every step.
\QED
\end{proof}

% ═════════════════════════════════════════════════════════════════════════════
\section{Model Identifiability}
% ═════════════════════════════════════════════════════════════════════════════

\subsection{Conditions for Expert Separability}

\begin{theorem}{Generic Identifiability of HNMC}{ident}
Suppose the following conditions hold:
\begin{enumerate}[label=\textbf{(C\arabic*)}, noitemsep]
  \item \textbf{Rank.}\; $P$ has full rank $K=7$.
  \item \textbf{Expert distinguishability.}\;
        For any $j\ne k$, there exists $(\bc,\bh)$ such that
        $p(\bc\mid S_j,\bh)\ne p(\bc\mid S_k,\bh)$.
  \item \textbf{Ergodicity.}\;
        The Markov chain induced by $P$ is irreducible and aperiodic.
\end{enumerate}
Then if two parameter sets $\lambda$ and $\lambda'$ yield the same joint
distribution, $p(\bc_{1:T}\mid\lambda)=p(\bc_{1:T}\mid\lambda')$ for all
$T$ and $\bc_{1:T}$, there exists a permutation $\sigma$ of $\{1,\ldots,K\}$
such that $\lambda' = \sigma(\lambda)$.
\end{theorem}

\begin{proof}
The proof adapts the tensor-decomposition approach of
Allman, Matias \& Rhodes~\cite{Allman2009} to neural emissions.

\medskip\noindent
\textbf{Step 1: Method of moments.}\;
Under (C3), the chain has a unique stationary distribution
$\pistar$ satisfying $P^\top\pistar=\pistar$.
The second-order observable moment satisfies
\[
  \E[\bc_t \bc_t^\top]
  = \sum_{k=1}^K \pistar_k \;\E\bigl[\by_t^{(k)}(\by_t^{(k)})^\top\bigr],
\]
a mixture decomposition indexed by $k$.

\medskip\noindent
\textbf{Step 2: Third-order tensor.}\;
Consider
$\mathbf{M}_3 = \E[\bc_{t-1}\otimes\bc_t\otimes\bc_{t+1}]$.
Under condition (C2), the expert emission means
$\{\bm\mu_k\}_{k=1}^K$ are linearly independent in $\R^d$ with $d\ge K$.
By \textbf{Kruskal's theorem}~\cite{Kruskal1977}:
a rank-$K$ tensor decomposition is unique (up to permutation of components)
when
\[
  \krank(A)+\krank(B)+\krank(C)\;\ge\; 2K+2,
\]
where $A,B,C$ are the factor matrices formed from the emission vectors.
Under (C1)--(C2), this inequality holds with probability~1 over generic
parameter draws (the set of parameter values violating it has Lebesgue
measure zero).

\medskip\noindent
\textbf{Step 3: Recovery of $P$.}\;
Given the identified emission parameters $\{\psi_k\}$ and $\pistar$,
the transition matrix $P$ is uniquely determined by
\[
  \E[\mathbf{e}_{z_t}\mathbf{e}_{z_{t+1}}^\top]
  = \diag(\pistar)\,P,
\]
which has full rank by (C1), so
$P = \diag(\pistar)^{-1}\,\E[\mathbf{e}_{z_t}\mathbf{e}_{z_{t+1}}^\top]$
is uniquely recovered.
Label permutation ambiguity is unavoidable in any mixture model and does not
affect predictive performance.  \QED
\end{proof}

\begin{remark}{Practical Identifiability on MELD}{meldident}
Condition~(C1) is satisfied by initialising $P$ from empirical MELD
inter-emotion transition counts (all 49 entries are positive in the corpus).
Condition~(C2) is monitored via the posterior entropy
$H(z_t\mid\bc_{1:t})$: if two experts collapse, an auxiliary expert-diversity
regulariser $\alpha\sum_{j\ne k}\mathrm{KL}(p_j\|p_k)^{-1}$ is added to the
objective.  Condition~(C3) holds because all seven emotions appear across
MELD's 1433 conversations.
\end{remark}

% ═════════════════════════════════════════════════════════════════════════════
\section{Convergence of Hybrid EM-Backprop Training}
% ═════════════════════════════════════════════════════════════════════════════

\subsection{Objective Function}

Given a dataset $\mathcal{D}=\{U^{(n)}\}_{n=1}^{\mathcal{N}}$ of
conversations, the training objective is the average marginal log-likelihood:
\begin{equation}
  \cL(\lambda)
  = \frac{1}{\mathcal{N}}
    \sum_{n=1}^{\mathcal{N}}
    \log p\!\bigl(\bc_{1:T_n}^{(n)} \mid \lambda\bigr)
  = \frac{1}{\mathcal{N}}
    \sum_{n=1}^{\mathcal{N}}
    \log \sum_{z_{1:T_n}} p\!\bigl(\bc_{1:T_n}^{(n)}, z_{1:T_n} \mid \lambda\bigr).
  \label{eq:obj}
\end{equation}

\begin{theorem}{Monotone Likelihood Improvement}{monotone}
The hybrid EM-backprop algorithm---alternating a closed-form M-step for $P$
(Baum-Welch) and a gradient M-step for $\{\psi_k\}$---satisfies
\[
  \cL(\lambda^{(s+1)}) \;\ge\; \cL(\lambda^{(s)})
  \quad\text{for all iterations }s.
\]
Consequently $\{\cL(\lambda^{(s)})\}$ is non-decreasing and bounded above,
hence convergent, and every limit point is a stationary point of $\cL$.
\end{theorem}

\begin{proof}
Define the EM auxiliary function
$\cQ(\lambda\mid\lambda^{(s)})
 = \E_{z\mid\bc,\,\lambda^{(s)}}[\log p(\bc,z\mid\lambda)]$.

\medskip\noindent
\textbf{EM lower bound.}\;
By Jensen's inequality applied to the concave log function,
\[
  \cL(\lambda) \;\ge\;
  \cQ(\lambda\mid\lambda^{(s)}) - H(z\mid\bc,\lambda^{(s)}),
\]
where the entropy $H$ does not depend on $\lambda$.  Therefore
\[
  \cL(\lambda^{(s+1)}) - \cL(\lambda^{(s)})
  \;\ge\;
  \cQ(\lambda^{(s+1)}\mid\lambda^{(s)}) - \cQ(\lambda^{(s)}\mid\lambda^{(s)}).
\]

\medskip\noindent
\textbf{$P$-step (closed form).}\;
The posterior occupation statistics are
\[
  \xi_{jk}(t) = p(z_{t-1}=j,\, z_t=k \mid \bc_{1:T},\,\lambda^{(s)}),
  \qquad
  \gamma_j(t) = \sum_k \xi_{jk}(t).
\]
The Baum-Welch update
\[
  P_{jk}^{\mathrm{new}}
  = \frac{\displaystyle\sum_t \xi_{jk}(t)}
         {\displaystyle\sum_t \gamma_j(t)}
\]
is the \emph{exact} maximiser of $\cQ$ over the simplex constraint on each
row of $P$, so
$\cQ(\lambda^{(s+1)}_P\mid\lambda^{(s)}) \ge \cQ(\lambda^{(s)}\mid\lambda^{(s)})$.

\medskip\noindent
\textbf{$\psi$-step (gradient).}\;
The Mamba SSM recurrence \cref{eq:ssm-h} involves compositions of matrix
exponentials, linear maps, and MLP projections---all smooth ($C^\infty$)
functions of $\psi_k$.  Therefore $\cQ$ is smooth in $\psi_k$, and its
gradient $\nabla_{\psi_k}\cQ$ is Lipschitz on any compact parameter set.
A single gradient ascent step with Armijo-Goldstein line search guarantees
$\cQ(\lambda_\psi^{(s+1)}\mid\lambda^{(s)}) \ge \cQ(\lambda^{(s)}_P
\mid\lambda^{(s)})$.

\medskip\noindent
\textbf{Combining.}\;
$\cL(\lambda^{(s+1)})
 \ge \cQ(\lambda^{(s+1)}\mid\lambda^{(s)}) - H
 \ge \cQ(\lambda^{(s)}\mid\lambda^{(s)}) - H
 = \cL(\lambda^{(s)})$.

\medskip\noindent
\textbf{Convergence.}\;
$\cL$ is bounded above for finite data.  The sequence is non-decreasing and
bounded, hence convergent.  By the generalised convergence theorem for EM
algorithms \cite{Wu1983}, every limit point satisfies
$\nabla_\lambda\cL(\lambda^*)=\mathbf{0}$.  \QED
\end{proof}

% ═════════════════════════════════════════════════════════════════════════════
\section{Strict Expressiveness over HMM and MoE}
% ═════════════════════════════════════════════════════════════════════════════

\subsection{Subsumption Results}

\begin{proposition}{HNMC Subsumes Classical HMM}{subhm}
Any $K$-state HMM with Gaussian emissions is a special case of HNMC.
\end{proposition}

\begin{proof}
Set $A_k=\mathbf{0}$, $C_k=\mathbf{0}$, $D_k=I$ in every expert.
Then $\bh_t^{(k)} = \Bbar_k(\bc_t)\,\bc_t \equiv \mathbf{0}$ (since
$A_k=\mathbf{0}$ implies $\Bbar_k=\Delta_k B_k$ and the state carries no
information forward), and
$\by_t^{(k)} = D_k\bc_t = \bc_t$.
The emission then reduces to
$p(\bc_t\mid S_k,\mathbf{0})
 = \mathcal{N}(\bc_t;\,\bm\mu_k,\Sigma_k)$
with $\bm\mu_k=b_k^{\mu}$ and fixed $\Sigma_k$---exactly a Gaussian HMM.
The transition matrix $P$ is unchanged.  \QED
\end{proof}

\begin{proposition}{HNMC Subsumes Standard MoE}{submoe}
Any $K$-expert Mixture-of-Experts model with input-gated routing is a
special case of HNMC.
\end{proposition}

\begin{proof}
Set $P = \frac{1}{K}\mathbf{1}_{K\times K}$ (uniform, memoryless transitions)
and initialise $\gamma_k(0)=\softmax(G\,\bc_1)_k$ for a learned gating matrix
$G$.  The forward recursion \cref{eq:fw-rec} then reduces to
\[
  \gamma_k(t)
  \propto p(\bc_t\mid S_k,\bh_{t-1}^{(k)})
         \cdot\frac{1}{K}\sum_j\gamma_j(t-1)
  = \frac{1}{K}\,p(\bc_t\mid S_k,\bh_{t-1}^{(k)}),
\]
which is equivalent to input-gated routing based solely on the current context.
Each expert $\psi_k$ retains its full Mamba SSM---this is exactly MoE with
Mamba experts and a memoryless router.  \QED
\end{proof}

\begin{theorem}{HNMC is Strictly More Expressive}{strict}
There exist distributions over emotion sequences realisable by HNMC but not
by \textup{(a)}~any HMM with parametric emissions, or
\textup{(b)}~any memoryless MoE.
\end{theorem}

\begin{proof}
\textbf{Part (a) -- against HMM.}\;
Consider a conversation in which the speaker remains in the \emph{anger}
regime ($z_t=\text{anger}$ for $t=1,\ldots,\tau$) with monotonically
increasing vocal pitch across utterances:
$z_t^{A,\text{pitch}} = t\cdot\delta$ for some $\delta>0$.
The predictive distribution $p(\bc_t\mid S_{\text{anger}})$ must depend on
$t$ within the regime.
A Gaussian HMM assigns
$p(\bc_t\mid S_k)=\mathcal{N}(\bm\mu_k,\Sigma_k)$
\emph{independently} of within-regime history, so it cannot represent this
trend.
The Mamba expert $\psi_{\text{anger}}$ accumulates the pitch trajectory in
its state $\bh_t^{(\text{anger})}$, and the emission mean
$W^\mu\bh_t^{(\text{anger})}$ tracks the monotone trend.  By \Cref{thm:ua},
the SSM approximates this functional arbitrarily well; no HMM can.

\medskip\noindent
\textbf{Part (b) -- against memoryless MoE.}\;
Let $P_{\text{anger},\text{anger}}=0.85$ (high persistence).
Consider two utterances $\bc$ and $\bc'$ with identical acoustic and
linguistic features, differing only in prior hidden state:
$z_{t-1}=\text{anger}$ (case~1) versus $z_{t-1}=\text{neutral}$ (case~2).
A memoryless router computes
$\gamma_k(t)=\softmax(G\,\bc_t)_k$---identical for both cases since
$\bc_t=\bc_{t}'$.
HNMC via \cref{eq:fw-rec} produces
\[
  \gamma_{\text{anger}}(t)
  \propto p(\bc\mid S_{\text{anger}},\bh)\cdot
         \sum_j P_{j,\,\text{anger}}\,\gamma_j(t-1),
\]
which differs between the two cases because $\gamma_j(t-1)$ differs.
Hence HNMC correctly distinguishes these cases; memoryless MoE cannot.  \QED
\end{proof}

% ═════════════════════════════════════════════════════════════════════════════
\section{MELD-Specific Theoretical Guarantees}
% ═════════════════════════════════════════════════════════════════════════════

\subsection{Structural Regularisation of Rare Classes}

\begin{proposition}{Geometric Decay of Rare-Class Activations}{rareclass}
Let $S_k$ be a rare emotion (e.g.\ \emph{fear} or \emph{disgust}) with
stationary probability $\pistar_k\ll 1$.
Define $P_{\max,k} = \max_{j\ne k} P_{jk}$.
Then, conditioned on no activation of $S_k$ for $\tau$ consecutive
utterances,
\[
  \E\bigl[\gamma_k(t) \mid z_{t-\tau}\ne k,\ldots, z_{t-1}\ne k\bigr]
  \;\le\;
  \pistar_k\cdot P_{\max,k}^{\,\tau}.
\]
\end{proposition}

\begin{proof}
Conditioned on $z_{t'}\ne k$ for $t'\in\{t-\tau,\ldots,t-1\}$, the forward
variable satisfies
\[
  \gamma_k(t)
  = p(\bc_t\mid S_k,\bh_{t-1}^{(k)})
    \sum_{j\ne k} P_{jk}\,\gamma_j(t-1)
  \le 1\cdot P_{\max,k}\sum_{j=1}^K\gamma_j(t-1)
  = P_{\max,k}\,\mathcal{Z}_t^{-1},
\]
where $\mathcal{Z}_t$ is the normalisation constant.
Applying the recursion backwards for $\tau$ steps and using
$\sum_j\gamma_j(0)=1$ with $\gamma_k(0)\le\pistar_k$, we obtain the
geometric bound $\pistar_k\cdot P_{\max,k}^{\,\tau}$.
Since $P_{\max,k}<1$ (irreducibility from (C3)), the bound decays
exponentially with the inter-activation gap $\tau$.
This provides a \emph{structural prior on rare emotions} stronger than
loss-reweighting alone, which only rescales the gradient without constraining
the temporal activation pattern.  \QED
\end{proof}

\subsection{Monotone Information Gain from Trimodal Fusion}

\begin{theorem}{Monotone Information Gain}{monotoneinfo}
Let $\cL_T$, $\cL_{TA}$, and $\cL_{TAV}$ denote the maximum marginal
log-likelihood achievable using text only, text$+$audio, and all three
modalities respectively.  Then
\[
  \cL_T \;\le\; \cL_{TA} \;\le\; \cL_{TAV},
\]
with strict inequalities whenever the additional modality carries emotion
information not already encoded in the previous ones.
\end{theorem}

\begin{proof}
The fused context $\bc_t=\MLP([\bz_t^T\|\bz_t^A\|\bz_t^V])$ in
\cref{eq:fusion} is parameterised by the MLP weights.  The text-only context
$\bc_t^T=\MLP([\bz_t^T\|\mathbf{0}\|\mathbf{0}])$ is the special case where
audio and video channels are zeroed, achievable by setting the corresponding
input-weight columns of the MLP to zero.

Since $\cL_{TAV}$ is maximised over \emph{all} MLP parameter values while
$\cL_T$ is maximised over the strict subset with zeroed audio/video weights,
we have $\cL_{TAV}\ge\cL_T$ by the feasibility argument.

\medskip\noindent
\textbf{Strict inequality.}\;
By the data-processing inequality for mutual information:
$I(z_t;\bc_t^{TAV})\ge I(z_t;\bc_t^T)$, with equality if and only if
$\bz_t^A$ and $\bz_t^V$ are conditionally independent of $z_t$ given
$\bz_t^T$.  On MELD, prosody (audio) and facial expression (video) provide
complementary emotion cues---e.g.\ sarcastic utterances have neutral
text but non-neutral prosody---so
$I(z_t;\bz_t^A\mid\bz_t^T)>0$ and
$I(z_t;\bz_t^V\mid\bz_t^T,\bz_t^A)>0$
for a non-trivial fraction of utterances, yielding strict inequalities.  \QED
\end{proof}

% ═════════════════════════════════════════════════════════════════════════════
\section{Main Theorem}
% ═════════════════════════════════════════════════════════════════════════════

\begin{theorem}{HNMC Theoretical Soundness for MELD}{main}
The Hidden Neural Markov Chain with $K=7$ Mamba experts and trimodal context
fusion satisfies all of the following:
\begin{enumerate}[label=\textbf{(\arabic*)}, leftmargin=2.2em]
  \item \textbf{[Validity]}\;
        It defines a proper probability measure over 7-class emotion sequences.
        \textup{(\Cref{thm:joint})}
  \item \textbf{[Tractability]}\;
        Exact marginal likelihoods are computable in $\mathcal{O}(49T)$ time
        via the forward algorithm.
        \textup{(\Cref{thm:forward})}
  \item \textbf{[Expressiveness]}\;
        Each expert universally approximates arbitrary continuous within-regime
        emotion dynamics.
        \textup{(\Cref{thm:ua})}
  \item \textbf{[Identifiability]}\;
        Model parameters are recoverable from data up to expert-label
        permutation under mild rank and ergodicity conditions.
        \textup{(\Cref{thm:ident})}
  \item \textbf{[Convergence]}\;
        Hybrid EM-backprop training monotonically improves the marginal
        log-likelihood and converges to a stationary point.
        \textup{(\Cref{thm:monotone})}
  \item \textbf{[Strict Subsumption]}\;
        HNMC strictly contains both classical HMMs and standard
        Mixture-of-Experts as degenerate special cases.
        \textup{(\Cref{thm:strict})}
  \item \textbf{[MELD Structure]}\;
        The Markov transition matrix provides automatic geometric regularisation
        of rare emotion classes and guarantees monotone information gain from
        trimodal fusion.
        \textup{(\Cref{prop:rareclass,thm:monotoneinfo})}
\end{enumerate}
Collectively, these results establish that HNMC is a theoretically
well-founded architecture with \emph{provable} advantages over its constituent
components for conversational multimodal emotion recognition on MELD.
\end{theorem}

\begin{proof}
Each numbered statement is the conclusion of the referenced theorem or
proposition, proved individually in \cref{sec:validity,sec:forward,sec:ua,sec:ident,sec:convergence,sec:expressiveness,sec:meld}.
The logical dependencies form a directed acyclic graph:
\Cref{lem:ssmstat} underlies \cref{thm:ua,thm:forward};
\cref{thm:ua} underlies \cref{thm:strict};
\cref{thm:joint} underlies all subsequent results.  \QED
\end{proof}

% ─────────────────────────────────────────────────────────────────────────────
% References
% ─────────────────────────────────────────────────────────────────────────────
\bigskip
\textcolor{RuleLine}{\rule{\linewidth}{0.4pt}}
\begin{thebibliography}{99}
\small

\bibitem{Allman2009}
Allman, E.~S., Matias, C., \& Rhodes, J.~A. (2009).
Identifiability of parameters in latent structure models with many observed
variables.
\textit{Annals of Statistics}, \textbf{37}(6A), 3099--3132.

\bibitem{Cybenko1989}
Cybenko, G. (1989).
Approximation by superpositions of a sigmoidal function.
\textit{Mathematics of Control, Signals, and Systems}, \textbf{2}(4),
303--314.

\bibitem{Gu2023}
Gu, A., \& Dao, T. (2023).
Mamba: Linear-time sequence modeling with selective state spaces.
\textit{arXiv}:2312.00752.

\bibitem{Hammer2000}
Hammer, B. (2000).
On the approximation capability of recurrent neural networks.
\textit{Neurocomputing}, \textbf{31}(1--4), 107--123.

\bibitem{Hornik1989}
Hornik, K., Stinchcombe, M., \& White, H. (1989).
Multilayer feedforward networks are universal approximators.
\textit{Neural Networks}, \textbf{2}(5), 359--366.

\bibitem{Kruskal1977}
Kruskal, J.~B. (1977).
Three-way arrays: Rank and uniqueness of trilinear decompositions.
\textit{Linear Algebra and its Applications}, \textbf{18}(2), 95--138.

\bibitem{Poria2018}
Poria, S., Hazarika, D., Majumder, N., Naik, G., Cambria, E., \& Mihalcea,
R. (2018).
MELD: A multimodal multi-relational emotion recognition dataset for
conversations.
\textit{arXiv}:1810.02508.

\bibitem{Rabiner1989}
Rabiner, L.~R. (1989).
A tutorial on hidden Markov models and selected applications in speech
recognition.
\textit{Proceedings of the IEEE}, \textbf{77}(2), 257--286.

\bibitem{Siegelmann1992}
Siegelmann, H.~T., \& Sontag, E.~D. (1992).
On the computational power of neural nets.
\textit{Proceedings of COLT 1992}, 440--449.

\bibitem{Wu1983}
Wu, C.~F.~J. (1983).
On the convergence properties of the EM algorithm.
\textit{Annals of Statistics}, \textbf{11}(1), 95--103.

\end{thebibliography}

\end{document}
